\documentclass[a4paper]{article}

\usepackage{amsmath}
\usepackage{amssymb}
\usepackage{stellar}
\usepackage{parskip}
\usepackage{fullpage}
\usepackage{enumitem}
\usepackage{wrapfig}
\usepackage{tikz}
\usepackage{witharrows}

\WithArrowsOptions{displaystyle}

\usetikzlibrary{arrows}
\usetikzlibrary{decorations.pathreplacing}
\usetikzlibrary{cd}

\title{Category theory}
\author{}
\date{}

\begin{document}

\maketitle
\tableofcontents

\pagebreak

\part{Fourier-Mukai transform}

\section{Abelian categories}

\begin{snippetdefinition}{Zero objct}
    Let \(\mathcal C\) be a category.
    An object \(0 \in \operatorname{Ob}(\mathcal C)\) is said to be a \emph{zero object}
    if it is both terminal and initial.
\end{snippetdefinition}

For any object \(X \in \operatorname{Ob}(\mathcal C)\),
\(\operatorname{Hom}_{\mathcal C}(0, X)\) and
\(\operatorname{Hom}_{\mathcal C}(X, 0)\) are both singletons.

\begin{snippetdefinition}{Zero morphism}
    Let \(\mathcal C\) be a category with a zero object \(0\).
    For any two objects \(X,Y \in \operatorname{Ob}(\mathcal C)\), the \emph{zero morphism}
    is defined as the unique morphism \(0_{XY} \colon X \fromto Y\) defined by
    \[
        0_{XY} \triangleq !_Y \circ !_X
    \]
    where \(!_Y \colon 0 \fromto Y\) is the unique morphism
    from \(0\) to \(Y\) and \(!_X \colon X \fromto 0\) is the unique morphism
    from \(X\) to \(0\).
\end{snippetdefinition}

\begin{snippetdefinition}{Category with zero morphisms}
    Let \(\mathcal C\) be a category.
    Then, \(\mathcal C\) is said to have \emph{zero morphisms} if there exists
    a fixed morphism \(0_{XY} \in \operatorname{Hom}_{\mathcal C}(X, Y)\) for every
    pair \(X,Y \in \operatorname{Ob}(\mathcal C)\) such that for any other
    object \(Z \in \operatorname{Ob}(\mathcal C)\) and morphisms \(f \colon X \fromto Y\) and \(g \colon Y \fromto Z\),
    the morphism \(0_{XY}\) absorbs composition, meaning that the following diagram commutes:
    \begin{center}
        % https://tikzcd.yichuanshen.de/#N4Igdg9gJgpgziAXAbVABwnAlgFyxMJZABgBpiBdUkANwEMAbAVxiRAA0QBfU9TXfIRQBGclVqMWbAJrdeIDNjwEiZYePrNWiELJ58lgoqPXVNUnQC1u4mFADm8IqABmAJwgBbJGRA4ISABMZpLaIC4g1Ax0AEYwDAAK-MpCIG5Y9gAWOHKuHt6Ion4BiADMIVps9rnh+T7U-khF5mHEAPrA7NJcNe5eQQ0l5RKVOu3A0pY9+rX9iL6NZRUWIOPsUzZcQA
        \begin{tikzcd}
            X \arrow[d, "f"'] \arrow[r, "0_{XY}"] \arrow[rd, "0_{XZ}"] & Y \arrow[d, "g"] \\
            Y \arrow[r, "0_{YZ}"]                                      & Z               
        \end{tikzcd}
    \end{center}
\end{snippetdefinition}

A category can also be defined to have zero morphisms without explicitly containing a zero object.
However, a category with a zero object always has zero morphisms.

\begin{snippetdefinition}{Biproduct}
    Let \(\mathcal C\) be a category with zero morphisms.
    Given a finite (possibly empty) collection of objects \(\{X_k\}_{k=1}^n\) in \(\operatorname{Ob}(\mathcal C)\),
    a \emph{biproduct} of this collection is an object \(\bigoplus_{k=1}^n X_k\)
    equipped with two families of morphisms:
    \begin{itemize}
        \item Projections: \(\pi_j \colon \bigoplus_{k=1}^n X_k \fromto X_j\) for each \(j \in \{1, \dots, n\}\)
        \item Injections: \(\iota_j \colon X_j \fromto \bigoplus_{k=1}^n X_k\) for each \(j \in \{1, \dots, n\}\)
    \end{itemize}
    such that:
    \begin{enumerate}
        \item \(\left(\bigoplus_{k=1}^n X_k, \pi_k\right)\) is a product for \(X_k\).
        \item \(\left(\bigoplus_{k=1}^n X_k, \iota_k\right)\) is a coproduct for \(X_k\).
        \item The morphisms satisfy the following relation for all \(j, k \in \{1, \dots, n\}\):
        \[
            \pi_j \circ \iota_k = \delta_{jk} = \begin{cases}
                \operatorname{id}_{X_j} & j=k \\
                0_{X_k X_j} & j \neq k
            \end{cases}
        \]
        where \(0_{X_k X_j}\) is the zero morphism.
    \end{enumerate}
    If the collection is empty, the biproduct is defined to be the zero object \(0\) of \(\mathcal C\).
\end{snippetdefinition}

In the category \(\mathbf{Set}\), the product and coproduct (cartesian product and disjoint union)
are essentially different. However, in abelian categories, such as the category of abelian groups
and vector spaces, the direct sum acts as both:
information can be extract by projecting on the singular spaces, or inserted by injecting from them.

\begin{snippetdefinition}{Preadditive category}
    A category \(\mathcal C\) is said to be \emph{preadditive} 
    if every hom-set \(\operatorname{Hom}_{\mathcal C}(X,Y)\) has the structure of an abelian group,
    and the composition of morphisms is bilinear.
    That is, for any morphisms \(f, f_1, f_2 \colon X \fromto Y\) and \(g, g_1, g_2 \colon Y \fromto Z\):
    \begin{align*}
        (g_1 + g_2) \circ f &= g_1 \circ f + g_2 \circ f \\
        g \circ (f_1 + f_2) &= g \circ f_1 + g \circ f_2
    \end{align*}
    The zero element of the abelian group \(\operatorname{Hom}_{\mathcal C}(X, Y)\) is the zero morphism \(0_{XY}\).
    This group structure allows us to freely use the addition \(f_1 + f_2\) and subtraction \(f_1 - f_2\) of any two parallel morphisms.
\end{snippetdefinition}

\begin{snippetdefinition}{Additive category}
    A category \(\mathcal C\) is said to be \emph{additive} if it is preadditive, has a zero object, and admits all finite biproducts.
\end{snippetdefinition}

In an additive category, the abelian group structure on the hom-sets is actually uniquely determined by the biproducts.
For any two parallel morphisms \(f, g \colon X \fromto Y\), their sum \(f + g\) can be explicitly defined as the composition:
\[
    f + g \triangleq \nabla_Y \circ (f \oplus g) \circ \Delta_X
\]
where:
\begin{itemize}
    \item \(\Delta_X \colon X \fromto X \oplus X\) is the \emph{diagonal morphism}, uniquely defined by the universal property of the product \(X \oplus X\) such that \(\pi_1 \circ \Delta_X = \operatorname{id}_X\) and \(\pi_2 \circ \Delta_X = \operatorname{id}_X\).
    \item \(f \oplus g \colon X \oplus X \fromto Y \oplus Y\) is the morphism naturally induced by \(f\) and \(g\) on the biproducts.
    \item \(\nabla_Y \colon Y \oplus Y \fromto Y\) is the \emph{codiagonal morphism}, uniquely defined by the universal property of the coproduct \(Y \oplus Y\) such that \(\nabla_Y \circ \iota_1 = \operatorname{id}_Y\) and \(\nabla_Y \circ \iota_2 = \operatorname{id}_Y\).
\end{itemize}
Because of this intrinsic structure, in abelian categories (which are a fortiori additive), we can seamlessly add, subtract, and evaluate linear combinations of morphisms. This structure is essential when studying chain complexes and homological algebra.

\begin{snippetdefinition}{Kernel}
    Let \(\mathcal C\) be a category with zero morphisms, and let \(f \colon X \fromto Y\) be a morphism in \(\mathcal C\).
    A \emph{kernel} of \(f\) is a pair \((K, \ker f)\), where \(K \in \operatorname{Ob}(\mathcal C)\) and \(\ker f \colon K \fromto X\), such that \(f \circ \ker f = 0_{KY}\), satisfying the following universal property:
    for every object \(Z \in \operatorname{Ob}(\mathcal C)\) and every morphism \(g \colon Z \fromto X\) such that \(f \circ g = 0_{ZY}\), there exists a unique morphism \(u \colon Z \fromto K\) such that the following diagram commutes:
    \begin{center}
        \begin{tikzcd}
            Z \arrow[rd, "g"] \arrow[d, "u", dashed] & \\
            K \arrow[r, "\ker f"'] & X \arrow[r, "f"] & Y
        \end{tikzcd}
    \end{center}
\end{snippetdefinition}

\begin{snippetdefinition}{Cokernel}
    Let \(\mathcal C\) be a category with zero morphisms, and let \(f \colon X \fromto Y\) be a morphism in \(\mathcal C\).
    A \emph{cokernel} of \(f\) is a pair \((C, \operatorname{coker} f)\), where \(C \in \operatorname{Ob}(\mathcal C)\) and \(\operatorname{coker} f \colon Y \fromto C\), such that \(\operatorname{coker} f \circ f = 0_{XC}\), satisfying the following universal property:
    for every object \(Z \in \operatorname{Ob}(\mathcal C)\) and every morphism \(g \colon Y \fromto Z\) such that \(g \circ f = 0_{XZ}\), there exists a unique morphism \(u \colon C \fromto Z\) such that the following diagram commutes:
    \begin{center}
        \begin{tikzcd}
            X \arrow[r, "f"] & Y \arrow[r, "\operatorname{coker} f"] \arrow[rd, "g"'] & C \arrow[d, "u", dashed] \\
            & & Z
        \end{tikzcd}
    \end{center}
\end{snippetdefinition}

The kernel is the equalizer of \((f, 0_{XY})\)
and measures how much information is lost via the map,
while the cokernel is the coequalizer of \((f, 0_{XY})\)
and measures how much of the domain is not covered.

\begin{snippetdefinition}{Image}
    Let \(\mathcal C\) be a category with zero morphisms, kernels, and cokernels.
    For any morphism \(f \colon X \fromto Y\) in \(\mathcal C\), the \emph{image} of \(f\), denoted \(\operatorname{im} f\), is defined as the kernel of its cokernel.
    Formally, it is the subobject of \(Y\) given by \(\operatorname{im} f \triangleq \ker(\operatorname{coker} f)\).
\end{snippetdefinition}

\begin{snippetdefinition}{Coimage}
    Let \(\mathcal C\) be a category with zero morphisms, kernels, and cokernels.
    For any morphism \(f \colon X \fromto Y\) in \(\mathcal C\), the \emph{coimage} of \(f\), denoted \(\operatorname{coim} f\), is defined as the cokernel of its kernel.
    Formally, it is the quotient object of \(X\) given by \(\operatorname{coim} f \triangleq \operatorname{coker}(\ker f)\).
\end{snippetdefinition}

\begin{snippetdefinition}{Abelian category}
    A category \(\mathcal A\) is said to be an \emph{abelian category} if it is an additive category satisfying the following two additional axioms:
    \begin{enumerate}
        \item \emph{Kernels and Cokernels}: Every morphism in \(\mathcal A\) has a kernel and a cokernel.
        \item \emph{Isomorphism Theorem}: For every morphism \(f\) in \(\mathcal A\), the canonical morphism \(\bar{f} \colon \operatorname{coim} f \fromto \operatorname{im} f\) is an isomorphism.
    \end{enumerate}
\end{snippetdefinition}

\section{Derived categories}

\begin{snippetdefinition}{Cochain Complex}
    Let \(\mathcal A\) be an abelian category. A \emph{cochain complex} \(C^\bullet\) in \(\mathcal A\) is a sequence of objects \(C^n \in \operatorname{Ob}(\mathcal A)\) for \(n \in \mathbb{Z}\), together with a sequence of morphisms \(d^n \colon C^n \fromto C^{n+1}\) called \emph{differentials}, such that the composition of any two consecutive differentials is the zero morphism:
    \[
        d^{n+1} \circ d^n = 0
    \]
    for all \(n \in \mathbb{Z}\).
    \begin{center}
        \begin{tikzcd}
            \cdots \arrow[r] & C^{n-1} \arrow[r, "d^{n-1}"] & C^n \arrow[r, "d^n"] & C^{n+1} \arrow[r, "d^{n+1}"] & \cdots
        \end{tikzcd}
    \end{center}
\end{snippetdefinition}

\begin{snippetdefinition}{Exact sequence}
    Let \(\mathcal A\) be an abelian category. A sequence of objects and morphisms
    \begin{center}
        \begin{tikzcd}
            \cdots \arrow[r] & C^{n-1} \arrow[r, "d^{n-1}"] & C^n \arrow[r, "d^n"] & C^{n+1} \arrow[r] & \cdots
        \end{tikzcd}
    \end{center}
    is said to be \emph{exact at \(C^n\)} if the image of the incoming morphism is equal to the kernel of the outgoing morphism:
    \[
        \operatorname{im} d^{n-1} = \ker d^n
    \]
    The entire sequence is called an \emph{exact sequence} if it is exact at every object \(C^n\).
\end{snippetdefinition}

Exactness means the elements that ``die'' in the second map are exactly those that came from the first map.
So nothing extra disappears, and nothing from the previous map is lost.

\begin{snippetdefinition}{Cohomology Object}
    Let \(C^\bullet\) be a cochain complex in an abelian category \(\mathcal A\).
    Because \(d^{n} \circ d^{n-1} = 0\), the image of \(d^{n-1}\) is naturally a subobject of the kernel of \(d^n\). 
    The \(n\)-th \emph{cohomology object} of \(C^\bullet\), denoted \(H^n(C^\bullet)\), is defined as the quotient object:
    \[
        H^n(C^\bullet) \triangleq \frac{\ker d^n}{\operatorname{im} d^{n-1}}
    \]
    which is well-defined since \(\mathcal A\) is an abelian category.
\end{snippetdefinition}

The condition \(d^{n} \circ d^{n-1} = 0\)
can be stated as \(\operatorname{im} d^{n-1} \subseteq \ker d^n\).
A sequence is said to be \emph{exact} if \(\operatorname{im} d^{n-1} \subseteq \ker d^n\).
The cohomology measures precisely the error towards being exact.
If \(H^n(C^\bullet) = 0\), the sequence is exact at that point.

\begin{snippetdefinition}{Morphism of Complexes (Chain Map)}
    Let \(\mathcal A\) be an abelian category, and let \((C^\bullet, d_C)\) and \((D^\bullet, d_D)\) be two cochain complexes in \(\mathcal A\).
    A \emph{morphism of complexes} (or \emph{chain map}) \(f^\bullet \colon C^\bullet \fromto D^\bullet\) is a sequence of morphisms \(f^n \colon C^n \fromto D^n\) in \(\mathcal A\) for each \(n \in \mathbb{Z}\), such that they commute with the differentials. 
    Explicitly, for all \(n \in \mathbb{Z}\), the following relation holds:
    \[
        d_D^n \circ f^n = f^{n+1} \circ d_C^n
    \]
    meaning that every square in the following diagram commutes:
    \begin{center}
        \begin{tikzcd}
            \cdots \arrow[r] & C^n \arrow[d, "f^n"'] \arrow[r, "d_C^n"] & C^{n+1} \arrow[d, "f^{n+1}"] \arrow[r] & \cdots \\
            \cdots \arrow[r] & D^n \arrow[r, "d_D^n"'] & D^{n+1} \arrow[r] & \cdots
        \end{tikzcd}
    \end{center}
\end{snippetdefinition}

\begin{snippetdefinition}{Quasi-isomorphism}
    Let \(f^\bullet \colon C^\bullet \fromto D^\bullet\) be a morphism of cochain complexes in an abelian category \(\mathcal A\).
    Because the sequence \(f^n\) commutes with the differentials, \(f^\bullet\) maps cycles to cycles and boundaries to boundaries, thus inducing a well-defined morphism on the cohomology objects for each \(n \in \mathbb{Z}\):
    \[
        H^n(f) \colon H^n(C^\bullet) \fromto H^n(D^\bullet)
    \]
    The morphism \(f^\bullet\) is called a \emph{quasi-isomorphism} (often abbreviated as \emph{quis}) if the induced morphism \(H^n(f)\) is an isomorphism in \(\mathcal A\) for every \(n \in \mathbb{Z}\).
\end{snippetdefinition}

If there exist a quasi-isomorphism between two
complexes, they contain the same (c)ohomological information.
The problem is that they are not necessarily actual isomorphisms.

\begin{snippetdefinition}{Category of Cochain Complexes}
    Let \(\mathcal A\) be an abelian category. The \emph{category of cochain complexes}, denoted \(\operatorname{Ch}(\mathcal A)\), is defined as follows:
    \begin{itemize}
        \item \emph{Objects}: Cochain complexes \((C^\bullet, d_C)\) in \(\mathcal A\).
        \item \emph{Morphisms}: Chain maps \(f^\bullet \colon C^\bullet \fromto D^\bullet\).
        \item \emph{Composition}: Given chain maps \(f^\bullet \colon C^\bullet \fromto D^\bullet\) and \(g^\bullet \colon D^\bullet \fromto E^\bullet\), their composition \((g \circ f)^\bullet\) is defined component-wise as \((g \circ f)^n \triangleq g^n \circ f^n\).
        \item \emph{Identity}: The identity morphism \(\operatorname{id}_{C^\bullet}\) consists of the identity morphisms \(\operatorname{id}_{C^n}\) on each object \(C^n\) for \(n \in \mathbb{Z}\).
    \end{itemize}
\end{snippetdefinition}

Since \(\mathcal A\) is abelian, \(\operatorname{Ch}(\mathcal A)\) is also an abelian category, where kernels and cokernels are computed component-wise.

In algebraic topology and homological algebra, two chain maps might not
be strictly identical, yet they might induce the exact same maps on the cohomology
objects. We thus need a way to formally relate them.
A \emph{chain homotopy} provides a continuous ``deformation'' between two chain maps,
similar to a topological homotopy between continuous functions.
It provides algebraic evidence that the difference between the two maps
is fundamentally constituted by boundaries of higher-dimensional elements,
and thus it vanishes when passing to the cohomology quotients.

\begin{snippetdefinition}{Chain Homotopy}
    Let \(f^\bullet, g^\bullet \colon C^\bullet \fromto D^\bullet\) be two chain maps between cochain complexes in an abelian category \(\mathcal A\).
    A \emph{chain homotopy} between \(f^\bullet\) and \(g^\bullet\) is a sequence of morphisms \(h^n \colon C^n \fromto D^{n-1}\) in \(\mathcal A\) for each \(n \in \mathbb{Z}\) such that:
    \[
        f^n - g^n = d_D^{n-1} \circ h^n + h^{n+1} \circ d_C^n
    \]
    If such a sequence \(\{h^n\}\) exists, \(f^\bullet\) and \(g^\bullet\) are said to be \emph{chain homotopic}, denoted by \(f^\bullet \sim g^\bullet\).
    
    \begin{center}
        % https://tikzcd.yichuanshen.de/#N4Igdg9gJgpgziAXAbVABwnAlgFyxMJZABgBpiBdUkANwEMAbAVxiRAB12BjKCHBAL6l0mXPkIoyARiq1GLNpx59BwkBmx4CRAMzlZ9Zq0QduvfiCEjN43aRnVDCk0vOrrY7SgBM+x-OMQAGEAPWAwAGopAUs1DU8JZCk-OSM2UMIrdVEtROSHVOcQABEQzLic2x97AwC2UvComIFZGCgAc3giUAAzACcIAFskMhAcCCQAViz+oanqcaQAFn80kygAfQzY3oHhxBWxicRvGb2kAHYF450zucQANmvL1aLN0sJqBjoAIxgGAAKlS8ID6WHaAAscDsQLN9skjkgHnd9pNno9XoEemUALTtMowuHLdFXEDfP6A4ESMkwHrQzFsbGNaJ4sKRaKE86INGIxAIpyBCEElHE3lPQqCtlNSwUARAA
        \begin{tikzcd}
            \cdots \arrow[r] & C^n \arrow[r, "d_C^n"] \arrow[d, "f^n-g^n"] \arrow[ld, "h^n"] & C^{n+1} \arrow[r] \arrow[d, "f^{n+1}-g^{n+1}"] \arrow[ld, "h^{n+1}"] & \cdots \\
            \cdots \arrow[r] & D^n \arrow[r, "d_D^n"']                                       & D^{n+1} \arrow[r]                                                    & \cdots
        \end{tikzcd}
    \end{center}
\end{snippetdefinition}

Notice that the subtraction \(f^n - g^n\) and the addition of compositions are well-defined operations precisely because \(\mathcal A\) is an additive category.

\begin{snippetdefinition}{Homotopy Category}
    Let \(\mathcal A\) be an abelian category. The \emph{homotopy category of cochain complexes}, denoted \(\mathbf{K}(\mathcal A)\), is defined as follows:
    \begin{itemize}
        \item \emph{Objects}: Cochain complexes \((C^\bullet, d_C)\) in \(\mathcal A\) (the same objects as \(\operatorname{Ch}(\mathcal A)\)).
        \item \emph{Morphisms}: Equivalence classes of chain maps modulo chain homotopy. Formally, the hom-sets are defined as:
        \[
            \operatorname{Hom}_{\mathbf{K}(\mathcal A)}(C^\bullet, D^\bullet) \triangleq \operatorname{Hom}_{\operatorname{Ch}(\mathcal A)}(C^\bullet, D^\bullet) / \sim
        \]
    \end{itemize}
\end{snippetdefinition}

\begin{snippetdefinition}{Localization of a Category}
    Let \(\mathcal C\) be a category and \(W\) be a class of morphisms in \(\mathcal C\). The \emph{localization} of \(\mathcal C\) with respect to \(W\), denoted \(\mathcal C[W^{-1}]\), is a category together with a functor \(Q \colon \mathcal C \fromto \mathcal C[W^{-1}]\) satisfying the following universal property:
    \begin{enumerate}
        \item For every morphism \(w \in W\), the image \(Q(w)\) is an isomorphism in \(\mathcal C[W^{-1}]\).
        \item For any other category \(\mathcal D\) and a functor \(F \colon \mathcal C \fromto \mathcal D\) such that \(F(w)\) is an isomorphism for all \(w \in W\), there exists a unique functor \(G \colon \mathcal C[W^{-1}] \fromto \mathcal D\) (up to natural isomorphism) such that the following diagram commutes:
        \begin{center}
            \begin{tikzcd}
                \mathcal C \arrow[rr, "F"] \arrow[rd, "Q"'] & & \mathcal D \\
                & \mathcal C[W^{-1}] \arrow[ru, "G"', dashed] &
            \end{tikzcd}
        \end{center}
    \end{enumerate}
\end{snippetdefinition}

Intuitively, the localization \(\mathcal C[W^{-1}]\) is obtained by formally adjoining inverses for all morphisms in \(W\), turning them into isomorphisms without introducing unnecessary equivalences.

\begin{snippetdefinition}{Derived Category}
    Let \(\mathcal A\) be an abelian category, and let \(\mathbf{K}(\mathcal A)\) be its homotopy category of cochain complexes. Let \(W\) be the class of all quasi-isomorphisms in \(\mathbf{K}(\mathcal A)\).
    The \emph{derived category} of \(\mathcal A\), denoted \(\mathbf{D}(\mathcal A)\), is defined as the localization of the homotopy category \(\mathbf{K}(\mathcal A)\) with respect to the class of quasi-isomorphisms \(W\):
    \[
        \mathbf{D}(\mathcal A) \triangleq \mathbf{K}(\mathcal A)[W^{-1}]
    \]
\end{snippetdefinition}

In the derived category, any quasi-isomorphism is formally inverted to become an actual equivalence. Thus, two complexes are isomorphic in \(\mathbf{D}(\mathcal A)\) if and only if they share the same homological information (i.e. they are connected by quasi-isomorphisms).

\section{Coherent sheaves}

\begin{snippetdefinition}{Presheaf on a Topological Space}
    Let \(X\) be a topological space, and let \(\mathcal{C}\) be a category. A \emph{presheaf} \(\mathcal{F}\) on \(X\) with values in \(\mathcal{C}\) consists of:
    \begin{enumerate}
        \item For every open set \(U \subseteq X\), an object \(\mathcal{F}(U)\) in \(\mathcal{C}\) (the elements are called \emph{sections} over \(U\)).
        \item For every inclusion of open sets \(V \subseteq U\), a morphism \(\operatorname{res}_{U,V} \colon \mathcal{F}(U) \fromto \mathcal{F}(V)\) in \(\mathcal{C}\) (called the \emph{restriction map}).
    \end{enumerate}
    This data must satisfy the following compatibility conditions:
    \begin{itemize}
        \item \(\mathcal{F}(\emptyset)\) is the terminal object of \(\mathcal{C}\).
        \item \(\operatorname{res}_{U,U} = \operatorname{id}_{\mathcal{F}(U)}\) for every open set \(U\).
        \item \(\operatorname{res}_{V,W} \circ \operatorname{res}_{U,V} = \operatorname{res}_{U,W}\) whenever \(W \subseteq V \subseteq U\).
    \end{itemize}
\end{snippetdefinition}

Intuitively, a presheaf assigns local algebraic or set-theoretic data to every open subset of a space, equipped with a consistent way to ``restrict'' that data to smaller interacting open subsets.

\begin{snippetdefinition}{Sheaf on a Topological Space}
    A presheaf \(\mathcal{F}\) with values in \(\mathbf{Set}\) on a topological space \(X\) is a \emph{sheaf} if it satisfies the following two additional conditions for any open set \(U \subseteq X\) and any open cover \(\{U_i\}_{i \in I}\) of \(U\):
    \begin{enumerate}
        \item \textbf{Locality}: If \(s, t \in \mathcal{F}(U)\) are sections such that \(\operatorname{res}_{U, U_i}(s) = \operatorname{res}_{U, U_i}(t)\) for all \(i \in I\), then \(s = t\).
        \item \textbf{Gluing}: Given a family of sections \(s_i \in \mathcal{F}(U_i)\) that agree on all overlaps, meaning \(\operatorname{res}_{U_i, U_i \intersection U_j}(s_i) = \operatorname{res}_{U_j, U_i \intersection U_j}(s_j)\) for all \(i, j \in I\), there exists a section \(s \in \mathcal{F}(U)\) such that \(\operatorname{res}_{U, U_i}(s) = s_i\) for all \(i \in I\).
    \end{enumerate}
\end{snippetdefinition}

The sheaf condition ensures that local data that matches on overlaps can always be uniquely glued together to form global structural data.

\begin{snippetdefinition}{Presheaf on a Category}
    Let \(\mathcal{C}\) be a small category. A \emph{presheaf} on \(\mathcal{C}\) with values in a category \(\mathcal{D}\) (typically \(\mathbf{Set}\)) is defined as a contravariant functor:
    \[
        P \colon \mathcal{C}^{\text{op}} \fromto \mathcal{D}
    \]
    The category of all presheaves on \(\mathcal{C}\) has these functors as objects and natural transformations between them as morphisms, taking the notation \([\mathcal{C}^{\text{op}}, \mathcal{D}]\) or \(\widehat{\mathcal{C}}\).
\end{snippetdefinition}

This categorical definition beautifully generalizes the topological one: if we regard the open sets of \(X\) as a poset category \(\mathcal{OP}(X)\) (where morphisms are inclusions), a topological presheaf is precisely a presheaf on the category \(\mathcal{OP}(X)\).

\begin{snippetdefinition}{Sheaf on a Category (with a Grothendieck Topology)}
    Let \(\mathcal{C}\) be a small category equipped with a Grothendieck topology \(J\) (making it a \emph{site}). A presheaf \(F \colon \mathcal{C}^{\text{op}} \fromto \mathbf{Set}\) is a \emph{sheaf} with respect to \(J\) if, for every object \(C \in \operatorname{Ob}(\mathcal{C})\) and every covering sieve \(S \in J(C)\), the natural map from the Yoneda embedding:
    \[
        \operatorname{Hom}_{\widehat{\mathcal{C}}}(\operatorname{Hom}_{\mathcal{C}}(-, C), F) \fromto \operatorname{Hom}_{\widehat{\mathcal{C}}}(S, F)
    \]
    is a bijection.
\end{snippetdefinition}

By substituting standard open covers with abstract covering sieves, Grothendieck topologies allow us to define sheaves entirely abstractly, unconditionally from an underlying point-set topology. Thus providing the foundation for topos theory.

\begin{snippetdefinition}{Ringed Space and Structure Sheaf}
    A \emph{ringed space} is a pair \((X, \mathcal{O}_X)\), where \(X\) is a topological space and \(\mathcal{O}_X\) is a sheaf of commutative rings on \(X\). The sheaf \(\mathcal{O}_X\) is called the \emph{structure sheaf} (or \emph{ringed sheaf}) of \(X\).
    
    A morphism of ringed spaces \((X, \mathcal{O}_X) \fromto (Y, \mathcal{O}_Y)\) is a uniquely mapped pair \((f, f^\sharp)\), where \(f \colon X \fromto Y\) is a continuous map and \(f^\sharp \colon \mathcal{O}_Y \fromto f_* \mathcal{O}_X\) is a morphism of sheaves of rings on \(Y\).
\end{snippetdefinition}

Ringed spaces associate a concrete ring of functions to every open set, essentially allowing us to study geometry purely through overarching algebraic structures.

\section{Fourier Transform}

\begin{snippetdefinition}{Topological Group}
    A \emph{topological group} is a group \(G\) endowed with a topology such that the group operations:
    \begin{itemize}
        \item Multiplication: \(\mu \colon G \cartesianprod G \fromto G\), \((x, y) \mapsto xy\)
        \item Inversion: \(\iota \colon G \fromto G\), \(x \mapsto x^{-1}\)
    \end{itemize}
    are continuous functions, where \(G \cartesianprod G\) is equipped with the product topology.
\end{snippetdefinition}

The continuity of these operations ensures that the algebraic structure of the group is compatible with its topological structure, allowing us to perform meaningful mathematical analysis on the group.

\begin{snippetdefinition}{Locally Compact Group}
    A topological group \(G\) is called a \emph{locally compact group} if its underlying topology is both locally compact and Hausdorff.
\end{snippetdefinition}

Locally compact Hausdorff groups provide the natural and primary setting for generalizing classical harmonic analysis, as local compactness guarantees the existence of a regular, non-trivial, translation-invariant measure.

\begin{snippetdefinition}{Haar Measure}
    Let \(G\) be a locally compact group. A \emph{left Haar measure} on \(G\) is a non-zero Radon measure \(\mu\) on the Borel \(\sigma\)-algebra of \(G\) that is left translation-invariant. That is, for every Borel set \(E \subseteq G\) and every \(g \in G\):
    \[
        \mu(gE) = \mu(E)
    \]
    where \(gE = \{gx \suchthat x \in E\}\). A right Haar measure is defined analogously.
\end{snippetdefinition}

The Haar measure provides a consistent way to integrate functions over a locally compact group, generalizing the translation-invariant Lebesgue measure on \(\realnumbers^n\). It is unique up to multiplication by a positive constant. For abelian groups, left and right Haar measures always coincide.

\begin{snippetdefinition}{Pontryagin Dual}
    Let \(G\) be a locally compact abelian (LCA) group. A \emph{character} of \(G\) is a continuous group homomorphism from \(G\) to the circle group \(\mathbb{T} = \{z \in \complexnumbers \suchthat |z| = 1\}\).
    The \emph{Pontryagin dual} of \(G\), denoted \(\widehat{G}\), is the set of all characters of \(G\):
    \[
        \widehat{G} = \operatorname{Hom}(G, \mathbb{T})
    \]
    equipped with point-wise multiplication and the compact-open topology.
\end{snippetdefinition}

Because \(G\) is a locally compact abelian group, its Pontryagin dual \(\widehat{G}\) is guaranteed to also be a locally compact abelian group. This dual group serves as the domain for the formal frequency spectrum in abstract harmonic analysis, much like how \(\realnumbers\) is its own dual.

\begin{snippetdefinition}{Abstract Fourier Transform}
    Let \(G\) be a locally compact abelian group with a chosen Haar measure \(\mu\). The \emph{Fourier transform} of a function \(f \in L^1(G, \mu)\) is the function \(\widehat{f} \colon \widehat{G} \fromto \complexnumbers\) defined by:
    \[
        \widehat{f}(\chi) = \int_G f(x) \overline{\chi(x)} \, d\mu(x)
    \]
    for every character \(\chi \in \widehat{G}\).
\end{snippetdefinition}

This generalized formulation of the abstract Fourier transform fundamentally encompasses the classical Fourier transform on \(\realnumbers\), Fourier series on \(\mathbb{T}\), and the Discrete Fourier Transform on finite cyclic groups, successfully unifying them under a single framework via Pontryagin duality.

\pagebreak

\part{Category theory exercises}

\section{Problems}

\begin{snippetexercise}{Exercise 8 from Chapter 1 of the book \emph{Sheaves in Geometry and Logic} by S. Mac Lane and I. Moerdijk}
    Consider a small category \(\mathcal C\). For each object \(B\)
    of \(\mathcal C\) there is a functor \(D_B \colon \mathcal C/B \fromto \mathcal C\) defined by taking the domain
    of each arrow to \(B\). Hence, each \(T \colon \mathcal C^{\text{op}} \fromto \mathbf{Set}\) yields
    \(T_B = T \circ D^{\text{op}} \colon {(\mathcal C/B)}^{\text{op}} \fromto \mathbf{Set}\).
    Define an exponential \(T^S\) by
    \[
        T^S(B) \triangleq \text{Hom}_{\widehat{\mathcal C/B}}(S_B, T_B),
    \]
    with the evident evaluations \(e_B \colon T^S(B) \cartesianprod S(B) \fromto T(B)\).
    Show that \(T^S\) with this evaluation \(e\) is indeed the exponential in the functor
    category \(\widehat{\mathcal C} = {\mathbf{Set}}^{\mathcal C^{\text{op}}}\).
\end{snippetexercise}

\begin{snippetexercise}{Exercise 7 from Chapter 2 of the book \emph{Sheaves in Geometry and Logic} by S. Mac Lane and I. Moerdijk}
    For any set \(T\), the constant presheaf \(T\) on a space \(X\) has
    \(T(U) = T\) for all open sets \(U\) in \(X\), with all restriction maps the identity.
    Show, using germs, that the associated ètale space is the projection \(p \colon X \cartesianprod T \fromto X\) of the
    product, where \(T\) has the discrete topology; conclude that the associated sheaf is te ``constant''
    sheaf \(\Delta_T\), for which \(\Delta_T(V)\) is the set of all locally constant functions
    \(V \fromto T\). Prove also that this defined a functor \(\Delta \colon \mathbf{Set} \fromto \mathbf{Sh}(X)\)
    which is left adjoint to the global sections functor \(\mathbf{Sh}(X) \fromto \mathbf{Set}\),
    \(F \colon \Gamma F(X)\).
\end{snippetexercise}

\begin{snippetexercise}{Exercise 8 from Chapter 3 of the book \emph{Sheaves in Geometry and Logic} by S. Mac Lane and I. Moerdijk}
    Let \(\mathcal C\) be a small category, and let \(P \colon \mathcal C^{\text{op}} \fromto \mathbf{Set}\)
    be a presheaf on \(\mathcal C\). Recall that the category
    \(\int_{\mathcal C} P\) of elements of \(P\) has as objects the pairs
    \((x, C)\) with \(x \in P(C)\), and as morphisms \((x, C) \fromto (x', C')\) the morphisms
    \(f \colon C \fromto C'\) in \(\mathcal C\) with the property that \(x' \cdot f = x\).
    \begin{enumerate}[label=\Alph*]
        \item Prove that there is an equivalence of categories
        \[
            \mathbf{Set}^{\mathcal C^{\text{op}}} /P \cong \mathbf{Set}^{{\left(\int_C P\right)}^{\text{op}}}
        \]
        \item Suppose that \(J\) is a topology on \(\mathcal C\), and that \(P \colon \mathcal C^{\text{op}} \fromto \mathbf{Set}\)
        is a \(J\)-sheaf. Describe a topology \(J'\) on \(\int_{\mathcal C} P\) such that
        the equivalence of \((a)\) restricts to an equivalence
        \[
            \mathbf{Sh}\left(
                \mathcal C, J
            \right) \cong \mathbf{Sh} \left(
                \int_C P, J'
            \right)
        \]
    \end{enumerate}
\end{snippetexercise}

\begin{snippetexercise}{}
    Let \(f\colon A \fromto B\) be an arrow in a Grothendieck topos \(\mathcal E\).
    Show that the pullback functor \(f^* \colon \mathcal E / B \fromto \mathcal E / A\) is
    faithful if and only if \(f\) is an epimorphism.
\end{snippetexercise}

\begin{snippetexercise}{}
    Recall that a point of a topos \(\mathcal E\)
    is a geometric morphism \(\mathbf{Set} \fromto \mathcal E\).
    Show that
    \begin{enumerate}[label=\Alph*]
        \item For any precorder \(\mathcal C\) and Grothendieck topology \(J\)
        on \(\mathcal C\), the points of the topos \(\mathbf{Sh}(\mathcal C, J)\) correspond
        precisely to the \emph{J-prime filters} on \(\mathcal C\) (by a \(J\)-prime filter
        on \(\mathcal C\) we mean a subset \(F\subseteq \mathcal C\) such that \(F\) is non-empty,
        \(a \in F\) implies \(b\in F\) whenever \(a\leq b\), for any \(a,b \in F\) there exists
        \(c\in F\) such that \(c \leq a\) and \(c\leq b\), and for any \(J\)-covering sieve
        \(\{a_i \fromto a \suchthat i \in I\}\) in \(\mathcal C\) if
        \(a \in F\) then there exists \(i\in I\) such that \(a_i \in F\)).
        \item For any frame \(L\), the points of the topos
        \(\mathbf{Sh}(L)\) correspond precisely to the frame homomorphisms
        \(L \fromto \{0,1\}\), equivalently to the \emph{completely prime filters}
        on \(L\) (i.e., the subsets \(S \subseteq L\) such that \(1 \in S, a \land b \in S\) if
        and only if \(a \in S\) and \(b \in S\), and for any family of elements
        \(\{a_i \suchthat i\in I\}\) whose join is \(a\), \(a \in S\) implies \(a_i \in S\) for
        some \(i\)).
        \item For any small category \(\mathcal C\) and any object \(c\) of \(\mathcal C\),
        \item there is a point \(ev_c \colon \mathbf{Set} \fromto [\mathcal C^{\text{op}}, \mathbf{Set}]\)
        of the topos \([\mathcal C^{\text{op}}, \mathbf{Set}]\) whose inverse image is the evaluation
        functor at the object \(c\).
    \end{enumerate}
\end{snippetexercise}

\section{Solutions}

\end{document}